Modern military operations increasingly depend on decisions made across cyberspace and the electromagnetic spectrum (EMS). An RF behavior observed before, during, or around a cyber operation can provide early evidence that an adversary is preparing to scan, interfere, replay, sustain access, or shift emissions. In the Quantifiable and Resilient Cyberspace Windows of Superiority (QR-CWoS) setting, this early evidence matters because the value of a response depends on timing: a defender must recognize useful precursor behavior before the opportunity to act has closed. An RF sensing layer for QR-CWoS therefore cannot merely classify familiar emissions. It must also preserve uncertainty when an observed signal behavior falls outside the trained taxonomy.

This requirement conflicts with the closed-world assumption used by most RF classifiers. A standard deep classifier is trained to assign each input to one of the known classes, and its softmax output provides no native mechanism for saying that an over-the-air (OTA) RF window does not belong to any known behavior. In contested RF environments, this assumption is unsafe. Jammers, spoofers, replay emitters, cooperative radios, protocol variants, and hardware/channel effects can all produce behaviors that differ from the available training set. The dangerous failure mode is not simply that the classifier is wrong; it is that the classifier can be wrong with high confidence, routing a genuinely novel behavior into the nearest known class exactly when the downstream system should be alerted to ambiguity.

Open-set recognition (OSR) addresses this failure mode by allowing a model to retain known-class recognition while rejecting samples that do not conform to the learned decision structure. Prior confidence-based OSR work, including varMax~\cite{baye2024varmax,broggi2025varmax}, shows that classifier outputs such as top-score gaps, logit variance, and energy-style evidence can expose unfamiliar inputs without retraining the backbone. Trott et al.~\cite{trott2026model} extend this idea to RF modulation classification and show that RF novelty can be strongly dependent on the composition of the known and unknown classes. Tiwari et al.~\cite{tiwaridqn} further show that softmax-derived confidence features can support a DQN-style known/unknown decision head in network intrusion detection. Together, these results suggest that an OSR decision layer can serve as a useful front end for larger cyber/EMS reasoning systems. They also expose a key constraint for our setting: the decision layer must control how many known RF behaviors it rejects while still preserving candidate unknowns for later analysis.

We study this problem for RF preliminary actions (PAs). A preliminary action is an observable RF behavior that may precede, enable, or contextualize a later cyber, EMS, or cyber-physical effect. The evaluated PA classes are Scan, Burst, Sustain, Hop, and Replay. These labels are intentionally not final ATT\&CK or EW technique labels. Instead, they are precursor-level evidence: Scan-like behavior can support reconnaissance hypotheses, Replay-like behavior can support compromise or unauthorized-command hypotheses, and frequency-agile or sustained behavior can support other EW-relevant interpretations. Known PA detections can therefore feed downstream attack-chain or response-planning modules, while rejected unknowns form a behavior pool for LLM-assisted label-making and continual RF behavior discovery.

This paper focuses on the sensing and triage layer required before those downstream modules can act. We introduce DQNGuard, a budgeted OSR decision layer for OTA RF preliminary-action recognition. DQNGuard operates on a multi-domain PA backbone that maps RF windows into IQ, FFT, DCT, and polar representations and produces closed-set logits, softmax probabilities, intermediate features, and a predicted PA. The decision layer then applies three stages: predicted-class conditional calibration, variance/energy guard evidence, and known-only budgeted thresholding. The threshold is selected from known calibration samples so that unknown detection is constrained by an explicit known-rejection budget rather than by an unconstrained search for novelty.

We evaluate DQNGuard on OTA RF captures spanning WiFi, Bluetooth, and Zigbee emissions and five PA behaviors. The experiments compare DQNGuard against varMax and a DQN-IDS-style confidence head, then analyze how surrogate-open calibration transfers across target unknown classes. The goal is not to claim a complete QR-CWoS response system or a final semantic labeling system for all novel emissions. Rather, the goal is to determine whether an RF front end can separate known preliminary-action evidence from nonconforming unknown behavior under controlled calibration assumptions.

Our contributions are as follows:
\begin{enumerate}
    \item We formulate open-world RF preliminary-action detection as a sensing layer for QR-CWoS, where known PA detections provide precursor evidence and rejected unknowns provide a pool for downstream label-making.
    \item We introduce DQNGuard, a budgeted OSR decision layer that combines predicted-class calibration, variance/energy guard evidence, and known-only thresholding to control known-class rejection while identifying unknown RF behaviors.
    \item We evaluate DQNGuard on an OTA preliminary-action dataset spanning three protocol families and five behavior classes under fixed-surrogate, Target--Surrogate Matrix, and target-open calibration regimes.
    \item We show that surrogate-open calibration is target-dependent, motivating future work on principled surrogate selection, pooled calibration, and continual RF behavior discovery.
\end{enumerate}

The remainder of the paper is organized as follows. Section~\ref{sec:related} summarizes related work in OSR, RF signal classification, and DQN-based intrusion detection. Section~\ref{sec:methodology} defines the PA detection problem, baseline OSR heads, and DQNGuard. Section~\ref{sec:experiments} describes the OTA dataset, training splits, calibration regimes, and metrics. Section~\ref{sec:results} reports the open-set detection results. Section~\ref{sec:discussion} discusses operational implications for QR-CWoS, ATT\&CK/EW precursor evidence, and continual label-making, and Section~\ref{sec:conclusion} concludes.
